\documentclass[12pt]{article}
\usepackage[utf8]{inputenc}
\usepackage[T1]{fontenc}
\usepackage{hyperref}
\title{Simba Language Resources for K-12 Curriculum Development}
\author{Compiled by Claude Sonnet 5.5 (claude-sonnet-5-5) for Andr\'as Kornai}
\date{}
\begin{document}
\maketitle

\section{Introduction}
Simba (ISO 639-3: sbw; also Himba, Himbaka, Ghehimba, Nsindak) is a Bantu language of Gabon (Guthrie B.302, Tsogo group), spoken in Ogooue-Lolo Province between Sindara and Mimongo. Wikipedia calls it moribund, with about 100 speakers (2007). Other sources give higher figures (3,000 speakers in 1990 according to the Endangered Languages Project). Wikitongues lists it as unwritten. Do not confuse it with the Himba dialect of Herero (Namibia) or with Zemba. No school use was found.

\section{Primary Educational Resources}
None found.

\section{Linguistic and Reference Materials}
\begin{itemize}
\item Rekanga, Jean-Paul (2000), \emph{Essai de grammaire himba (langue bantoue du Gabon, B36)}, doctoral thesis, Universit\'e Libre de Bruxelles, two volumes (Glottolog record). \url{https://glottolog.org/resource/reference/id/92848}
\item Wikipedia, Himba language. \url{https://en.wikipedia.org/wiki/Himba_language}
\item ASJP Simba wordlist. \url{https://asjp.clld.org/languages/SIMBA}
\item Endangered Languages Project record. \url{https://endangeredlanguages.com/lang/4616}
\end{itemize}
No dictionary was found.

\section{Religious and Community Resources}
None found.

\section{Digital and Audio Resources}
Wikitongues has a language page for sbw (alternate names, status). Audio or video there was not confirmed. \url{https://wikitongues.org/languages/sbw/}

\section{Curriculum Development Context}
French is the language of instruction in Gabon. Simba is unwritten, and the speaker base is small and old. Rekanga's thesis, in French, is the main description. Any work would be documentation and revitalisation first, with no existing teaching material to build on.
\end{document}